\chapter{الكسور: الخطوات الأولى}\label{ch:g6:fractions}

تقاسم ثلاث بيتزات بالعدل بين أربعة أشخاص يعطي كل واحد\,\dots\
لا عددًا صحيحًا من البيتزات. الكسور هي الأعداد التي اُخترعت
للتقاسم. ويبني هذا الفصل الصورة: \omterm{def:g6:fractions:def}{الكسر} جزءًا من كلّ، ونقطةً على
مستقيم الأعداد، وناتجًا مضبوطًا لقسمة.

\section{ماذا يعني الكسر}

\begin{definition}[الكسر]\label{def:g6:fractions:def}
اِقسم الوحدة إلى $b$ أجزاء متساوية، وخذ $a$ منها: المقدار
الناتج هو \emph{الكسر}\index{كسر}
\[
\frac{a}{b}
\qquad
\begin{array}{l}
a \text{: \emph{البسط} --- عدد الأجزاء التي نأخذها؛}\\
b \text{: \emph{المقام} --- إلى كم جزءًا قُسّمت الوحدة.}
\end{array}
\]
\end{definition}

\begin{omfigure}
\begin{tikzpicture}[scale=0.55]
  \foreach \i in {0,...,3} \draw[omExo] (\i*2,0) rectangle (\i*2+2,1.2);
  \foreach \i in {0,...,2} \fill[omDef!40] (\i*2+0.05,0.05)
    rectangle (\i*2+1.95,1.15);
  \draw[thick] (0,0) rectangle (8,1.2);
  \node[below, font=\small] at (4,-0.15) {$\frac34$ من اللوح: اِقسم إلى
  $4$، وخذ $3$};
  \begin{scope}[xshift=11.5cm]
  \foreach \i in {0,...,7} \draw[omExo] (\i,0) rectangle (\i+1,1.2);
  \foreach \i in {0,...,4} \fill[omThm!35] (\i+0.05,0.05)
    rectangle (\i+0.95,1.15);
  \draw[thick] (0,0) rectangle (8,1.2);
  \node[below, font=\small] at (4,-0.15) {$\frac58$ من اللوح: اِقسم إلى
  $8$، وخذ $5$};
  \end{scope}
\end{tikzpicture}

{\small كسران من اللوح نفسه. \omterm{def:g4:fractions:def}{المقام} يبيّن دقّة التقطيع،
\omterm{def:g4:fractions:def}{والبسط} يبيّن كم نأخذ منه.}
\end{omfigure}

\begin{example}\label{ex:g6:fractions:morethanone}
يمكن أن يكون \omterm{def:g6:fractions:def}{الكسر} أكبر من $1$: \omterm{def:g6:fractions:def}{فالكسر} $\frac74$ يعني سبعة أرباع ---
وحدة كاملة (أربعة أرباع) وثلاثة أرباع أخرى:
\[
\frac74 = 1 + \frac34 .
\]
\end{example}

\begin{proposition}[الكسور على مستقيم الأعداد]\label{prop:g6:fractions:line}
لوضع $\frac ab$ على مستقيم الأعداد، اِقسم كل فترة واحدية إلى $b$
أجزاء متساوية، وعُدّ $a$ جزءًا انطلاقًا من $0$.
\end{proposition}

\begin{omfigure}
\begin{tikzpicture}[scale=1.5]
  \draw[-Stealth] (-0.2,0) -- (7.6,0);
  \foreach \i in {0,4,8,12,16,20,24,28}
    \draw ({\i*0.25},-0.09) -- ({\i*0.25},0.09);
  \foreach \i in {0,...,28} \draw ({\i*0.25},-0.05) -- ({\i*0.25},0.05);
  \node[below, font=\small] at (0,-0.1) {$0$};
  \node[below, font=\small] at (1,-0.1) {$1$};
  \node[below, font=\small] at (2,-0.1) {$2$};
  \node[below, font=\small] at (3,-0.1) {$3$};
  \node[below, font=\small] at (4,-0.1) {$4$};
  \node[below, font=\small] at (5,-0.1) {$5$};
  \node[below, font=\small] at (6,-0.1) {$6$};
  \node[below, font=\small] at (7,-0.1) {$7$};
  \fill[omDef] (0.75,0) circle (1.6pt)
    node[above, font=\small] {$\frac34$};
  \fill[omThm] (1.75,0) circle (1.6pt)
    node[above, font=\small] {$\frac74$};
  \fill[omMeth] (5.5,0) circle (1.6pt)
    node[above, font=\small] {$\frac{22}{4}$};
\end{tikzpicture}

{\small المستقيم مدرّج بالأرباع: $\frac34$ قبل $1$،
و $\frac74$ بين $1$ و $2$، و $\frac{22}{4} = 5 + \frac24$ في
منتصف ما بين $5$ و $6$.}
\end{omfigure}

\begin{theorem}[الكسر قسمة]\label{thm:g6:fractions:quotient}
\omterm{def:g6:fractions:def}{الكسر} $\frac ab$ هو العدد الذي إذا ضُرب في $b$ أعطى
$a$:
\[
\frac ab \times b = a .
\]
وبعبارة أخرى، $\frac ab$ هو \omterm{def:g3:division:remainder}{خارج القسمة} المضبوط للعدد $a$ على $b$ --- حتى
حين لا «تنتهي القسمة بالضبط».
\end{theorem}
\admitted

\begin{example}\label{ex:g6:fractions:pizzas}
ثلاث بيتزات لأربعة أشخاص: نصيب كل واحد $\frac34$ من بيتزا.
والتحقّق بالمبرهنة: $4$ أنصبة من $\frac34$ تعطي
$\frac34 \times 4 = 3$ بيتزات. وبعض الكسور \omterm{def:g6:decimals:places}{أعداد عشرية}
($\frac34 = 0.75$)، وبعضها ليس كذلك: $\frac13 = 0.333\dots$ لا يتوقف أبدًا.
\omterm{def:g6:fractions:def}{والكسر} هو القيمة \emph{المضبوطة}؛ أما $0.33$ \omterm{def:g6:decimals:rounding}{فتقريب} لا غير.
\end{example}

\section{الكسور المتساوية}

\begin{proposition}[كسور متكافئة]\label{prop:g6:fractions:equal}
لا يتغيّر \omterm{def:g6:fractions:def}{الكسر} إذا ضُرب بسطه ومقامه معًا (أو قُسما معًا) على العدد
نفسه غير المعدوم:
\[
\frac{a}{b} = \frac{a \times k}{b \times k} .
\]
\end{proposition}

\begin{proof}[فكرة البرهان]
تقطيع كل جزء من التقاسم إلى $k$ قطعًا أصغر يضرب في $k$ عددَ القطع
المأخوذة وعددَ القطع كلها معًا، لكن المقدار المأخوذ لا يتغيّر.
والصورة أبلغ:
\end{proof}

\begin{omfigure}
\begin{tikzpicture}[scale=0.6]
  \foreach \i in {0,1,2} \draw[omExo] (\i*2.4,0) rectangle (\i*2.4+2.4,1.2);
  \foreach \i in {0,1} \fill[omDef!40] (\i*2.4+0.05,0.05)
    rectangle (\i*2.4+2.35,1.15);
  \draw[thick] (0,0) rectangle (7.2,1.2);
  \node[below, font=\small] at (3.6,-0.15) {$\frac23$};
  \begin{scope}[xshift=10.5cm]
  \foreach \i in {0,...,5} \draw[omExo] (\i*1.2,0) rectangle (\i*1.2+1.2,1.2);
  \foreach \i in {0,...,3} \fill[omDef!40] (\i*1.2+0.05,0.05)
    rectangle (\i*1.2+1.15,1.15);
  \draw[thick] (0,0) rectangle (7.2,1.2);
  \node[below, font=\small] at (3.6,-0.15) {$\frac46$};
  \end{scope}
\end{tikzpicture}

{\small $\frac23 = \frac46$: تقطيع كل ثلث إلى قسمين يضاعف
العددين، في الأعلى وفي الأسفل، دون أن يغيّر \omterm{def:g4:measure:area}{المساحة} المظللة.}
\end{omfigure}

\begin{method}[اِختزال كسر]\label{met:g6:fractions:simplify}
اِبحث عن عدد يقسم \omterm{def:g4:fractions:def}{البسط} \emph{و} \omterm{def:g4:fractions:def}{المقام} معًا، واقسمهما عليه.
وكرّر حتى لا يبقى قاسم مشترك:
\[
\frac{12}{18} = \frac{12 \div 2}{18 \div 2} = \frac{6}{9}
= \frac{6 \div 3}{9 \div 3} = \frac{2}{3}.
\]
\end{method}

\section{أخذ كسر من مقدار}

\begin{method}[كسر من مقدار]\label{met:g6:fractions:ofquantity}
لحساب $\frac ab$ من مقدار، اِقسم المقدار على $b$ (فتحصل على جزء
واحد)، ثم اِضرب في $a$ (وهو عدد الأجزاء):
\[
\frac ab \text{ من } Q = (Q \div b) \times a .
\]
(والضرب في $a$ أولًا ثم القسمة على $b$ يعطي النتيجة نفسها ---
فاختر الترتيب الذي يجعل الأعداد أسهل.)
\end{method}

\begin{example}\label{ex:g6:fractions:ofquantity}
اِحسب $\frac35$ من $40$ يورو، خطوة خطوة:
\begin{enumerate}
  \item خُمس العدد $40$: $40 \div 5 = 8$؛
  \item ثلاثة أخماس: $8 \times 3 = 24$.
\end{enumerate}
إذن $\frac35$ من $40$ يورو هو $24$ يورو. والتحقّق: الباقي
$\frac25$ يساوي $16$ يورو، و $24 + 16 = 40$.
\end{example}

\section{تمارين}

\begin{exercise}[$\star$]\label{exo:g6:fractions:1}
لكل وصف صورة، اُكتب الكسر: كعكة قُطّعت إلى $8$
شرائح، أُكل منها $3$؛ ولوح شوكولاتة من $12$ مربعًا، أُكل منه $7$؛ وفطيرة
قُطّعت إلى $6$، أُكل منها $6$.
\end{exercise}

\begin{exercise}[$\star$]\label{exo:g6:fractions:2}
اُرسم لوحًا مقسومًا إلى $5$ أجزاء متساوية ولوّن منه $\frac45$. هل
$\frac45$ أصغر من $1$ أم أكبر؟ وماذا عن $\frac65$؟
\end{exercise}

\begin{exercise}[$\star$]\label{exo:g6:fractions:3}
ضع على مستقيم أعداد مدرّج بالأثلاث: $\frac13$؛ \ $\frac53$؛
\ $2$؛ \ $\frac73$. أيّ عددين من هذه على المسافة نفسها من
$2$؟
\end{exercise}

\begin{exercise}[$\star$]\label{exo:g6:fractions:4}
اُكتب كل كسر على صورة عدد صحيح زائد كسر أصغر من $1$
(كما في $\frac74 = 1 + \frac34$):
\[
\frac{9}{4}, \qquad \frac{13}{5}, \qquad \frac{12}{6}, \qquad
\frac{25}{8} .
\]
\end{exercise}

\begin{exercise}[$\star$]\label{exo:g6:fractions:5}
أكمل حتى تتساوى الكسور:
\[
\frac{2}{3} = \frac{?}{12}, \qquad
\frac{15}{20} = \frac{3}{?}, \qquad
\frac{7}{10} = \frac{21}{?}, \qquad
\frac{?}{6} = \frac{10}{12}.
\]
\end{exercise}

\begin{exercise}[$\star$]\label{exo:g6:fractions:6}
اِختزل إلى أبسط صورة ممكنة: $\dfrac{6}{8}$؛ \quad $\dfrac{15}{25}$؛
\quad $\dfrac{18}{24}$؛ \quad $\dfrac{35}{7}$.
\end{exercise}

\begin{exercise}[$\star$]\label{exo:g6:fractions:7}
اِحسب: $\frac12$ من $86$؛ \quad $\frac34$ من $60$؛ \quad $\frac25$ من
$35$؛ \quad $\frac{7}{10}$ من $250$.
\end{exercise}

\begin{exercise}[$\star$]\label{exo:g6:fractions:8}
أيّهما أكبر: $\frac12$ أم $\frac13$ (من الكعكة نفسها)؟ و $\frac25$
أم $\frac35$؟ و $\frac34$ أم $\frac78$ (اُكتبهما بالمقام
$8$)؟
\end{exercise}

\begin{exercise}[$\star\star$]\label{exo:g6:fractions:9}
أعطِ القيمة العشرية للكسور التي لها قيمة عشرية، وقل أيّ
كسر ليست له: $\frac12$؛ \ $\frac34$؛ \ $\frac13$؛ \ $\frac{7}{10}$؛
\ $\frac{9}{4}$.
\end{exercise}

\begin{exercise}[$\star\star$]\label{exo:g6:fractions:10}
في قسم $28$ تلميذًا. ثلاثة أرباعهم يأتون إلى المدرسة
بالحافلة، والباقون يمشون. فكم تلميذًا يمشي؟ حُلّ خطوة خطوة، وتحقّق
من أن مجموع الفئتين يساوي $28$.
\end{exercise}

\begin{exercise}[$\star\star$]\label{exo:g6:fractions:11}
أنفق لؤي $\frac23$ من مدّخراته على لعبة ثمنها $18$ يورو. فكم
كان معه من المال قبل ذلك؟ (ثلث مدّخراته هو نصف
$18$\,\dots\ فكّر لماذا، أو اِبحث أولًا عن قيمة جزء واحد.)
\end{exercise}

\begin{exercise}[$\star\star\star$]\label{exo:g6:fractions:12}
خزّان مملوء إلى $\frac58$ من سعته. وبعد استعمال $10$ لترات
صار نصف ممتلئ. فما سعة الخزّان؟ (أيّ كسر من
الخزّان تمثّل $10$ لترات؟)
\end{exercise}

\section{مسألة: لوح الشوكولاتة الذي لا تنتهي منه أبدًا}

\begin{problem}[{مسألة نهاية الأسبوع --- أنصاف الأنصاف: المجموع
$\frac12 + \frac14 + \frac18 + \dots$ يزحف نحو $1$ دون أن
يبلغه أبدًا}]
\label{pb:g6:fractions:1}
خذ لوح شوكولاتة. كُلْ نصفه. ثم كُلْ نصف \emph{ما
بقي}. ثم نصف ما بقي مرة أخرى، وهكذا. ويبدو أمران صحيحين
في آن واحد: أنك لا تنهي اللوح أبدًا (فشيء ما يبقى دائمًا)، ومع
ذلك تأكل منه كل شيء تقريبًا. وتحوّل هذه المسألة الشعورين إلى
عبارتين مضبوطتين عن الكسور --- وتلتقي في الطريق عدّاءً من
اليونان القديمة، ولغز تقاسم يُحلّ بسكين لا يقطع إلا
أنصافًا.

\medskip
\noindent\textbf{الجزء الأول --- أنصاف الأنصاف.}
اللوح مستطيل من $16$ مربعًا متساويًا.
\begin{enumerate}
  \item القضمة الأولى: تأكل نصف اللوح. فكم مربعًا يساوي
        ذلك؟ اُكتب القضمة على صورة كسر من اللوح بطريقتين،
        بالمقام $2$ ثم بالمقام $16$
        (\cref{prop:g6:fractions:equal}).
  \item القضمة الثانية: نصف ما بقي. فكم مربعًا؟ بيّن،
        بالصورة أو بالكسور المتكافئة، أن هذه القضمة
        هي $\frac14$ من اللوح كله: \emph{فنصف النصف
        رُبع}.
  \item القضمتان الثالثة والرابعة، بالقاعدة نفسها: أعطِ كلًّا منهما
        بالمربعات وبكسر من اللوح. وماذا يحدث
        للمقام من قضمة إلى التي بعدها؟
  \item بعد القضمات الأربع، كم مربعًا أُكل في
        المجموع؟ وأيّ كسر من اللوح هو ذلك، وأيّ كسر
        بقي؟
  \item يقول السؤال 4، مكتوبًا بالأجزاء من ستة عشر:
        \[
        \frac{8}{16} + \frac{4}{16} + \frac{2}{16} + \frac{1}{16}
        = \frac{15}{16},
        \qquad\text{أي}\qquad
        \frac12 + \frac14 + \frac18 + \frac{1}{16}
        = \frac{15}{16} .
        \]
        تحقّق من كل واحدة من عمليات إعادة الكتابة الأربع
        ($\frac12 = \frac{8}{16}$، وهكذا).
\end{enumerate}

\medskip
\noindent\textbf{الجزء الثاني --- السُّلَّم نحو $1$.}
لمواصلة القضم، تخيّل لوحًا أدقّ من $64$ مربعًا.
\begin{enumerate}[resume]
  \item اُذكر أحجام القضمات الست الأولى، بالمربعات
        ($32$، ثم\,\dots)، وتحقّق من أنه بعد القضمة السادسة
        يبقى مربع واحد بالضبط. وأيّ كسر من اللوح
        يبقى؟
  \item بعد كل قضمة، أيّ كسر من اللوح يبقى؟ اُكتب
        القائمة ($\frac12$، $\frac14$، \dots) إلى القضمة
        السادسة، وصِف النمط. ولو كان سكين خيالي
        يستطيع أن يشطر حتى المربع الأخير، فماذا يبقى بعد
        قضمة سابعة؟
  \item اِشرح لماذا لا يمكن للمجموع المأكول أن يبلغ $1$
        \emph{أبدًا}، مهما كثُر عدد القضمات. (قارن مع
        التسعات العشرين في \cref{pb:g6:decimals:1}: أهي القصة نفسها؟)
  \item ومع ذلك يتجاوز المجموع أيّ هدف دون $1$. يتحدّاك
        صديق أن تأكل أكثر من $\frac{99}{100}$ من
        اللوح. فأيّهما أصغر، $\frac{1}{128}$ أم
        $\frac{1}{100}$ --- وبعد أيّ قضمة يُربح التحدي؟
  \item على مستقيم أعداد من $0$ إلى $1$ مدرّج بالأجزاء من ستة عشر،
        ضع المجاميع المأكولة بعد كل قضمة من القضمات الأربع
        الأولى: $\frac12$، $\frac34$، $\frac78$، $\frac{15}{16}$.
        وأين يقع كل مجموع جديد، بالمقارنة مع
        السابق ومع $1$؟
\end{enumerate}

\medskip
\noindent\textbf{الجزء الثالث --- عدّاء زينون، وسكين
لا يشطر إلا إلى نصفين.}
\begin{enumerate}[resume]
  \item قبل أربعة وعشرين قرنًا، روى الفيلسوف اليوناني زينون
        هذه الحكاية: لكي يبلغ العدّاء جدارًا، عليه أولًا أن
        يقطع نصف المسافة، ثم نصف ما بقي، ثم
        نصف ما بقي مرة أخرى\,\dots\ «فالعدّاء لا يبلغ
        الجدار أبدًا!» باستعمال السؤالين 8 و9، قل ما هو
        \emph{الصحيح} في حكاية زينون، وأين الفخّ.
        (أيحتاج العدّاء حقًّا إلى زمن لا ينتهي ليقطع كل تلك
        الخطوات؟)
  \item لغز التقاسم: اِقسم $15$ لوح شوكولاتة متطابقًا
        بالعدل بين $16$ طفلًا، مستعملًا القطع إلى أنصاف فقط
        (لوحٍ، أو نصف لوح، أو ربع لوح\,\dots).
        صِف كيف تقطع: كم لوحًا يُقطع إلى أنصاف،
        وكم إلى أرباع، وإلى أثمان، وإلى أجزاء من ستة عشر،
        بحيث يأخذ كل طفل قطعة واحدة من كل حجم.
  \item تحقّق من تقاسمك: اُكتب نصيب كل طفل على صورة مجموع
        كسور، واحسبه بالأجزاء من ستة عشر، وتأكّد من أن
        الأنصبة $16$ كلها تستهلك الألواح $15$ بالضبط.
  \item في القسم المجاور، تُقسم $7$ ألواح بالعدل بين
        $8$ أطفال. فأيّ طفل يأخذ شوكولاتة أكثر --- طفل
        من المجموعة ذات $16$ أم طفل من المجموعة ذات $8$؟ (اُكتب
        النصيبين بالمقام $16$، \cref{exo:g6:fractions:8}.)
  \item الخاتمة: يزعم صديق أنه بعدد \emph{لا نهائي}
        من القضمات يأكل المرء اللوح \emph{كله} بالضبط.
        اُرسم مربعًا وظلّل القضمات داخله: نصف
        المربع، ثم رُبعًا، ثم ثُمنًا\,\dots\ فماذا
        تلاحظ على الركن غير المظلل؟ اُذكر بعناية
        الحقيقتين الصحيحتين في هذه المسألة كلها: ما الذي يبقى
        بعد كل عدد منته من القضمات، وكم يصغر ذلك
        \omterm{def:g3:division:remainder}{الباقي}. (وأما ما يعنيه «\omterm{def:g1:addition:def}{المجموع} اللانهائي» بالضبط
        فقصة تُروى في مجلّد المرحلة الثانوية.)
\end{enumerate}
\end{problem}
